\[ %%%%%%%%%%%%%%%%%%%%%%%%%%%%%%%%%%%%%%%%%%%%%%%%%%%%%%%%%%%%%%%%%%%%%%%%%%%%%%%% \subsection{Neural Operators for Learning Dynamical Systems} \label{sec:neural_operators} %%%%%%%%%%%%%%%%%%%%%%%%%%%%%%%%%%%%%%%%%%%%%%%%%%%%%%%%%%%%%%%%%%%%%%%%%%%%%%%% Neural operator learning has emerged as a powerful paradigm for accelerating scientific simulations by learning mappings between infinite-dimensional function spaces rather than approximating individual solution instances. Unlike conventional neural networks that operate on fixed-dimensional input-output vectors, neural operators aim to approximate solution operators of physical systems, enabling rapid prediction across varying initial conditions, boundary conditions, and physical parameters. A general formulation of an operator learning problem can be expressed as: \begin{equation} \mathcal{G}: u \rightarrow v, \end{equation} where $\mathcal{G}$ represents an unknown nonlinear operator mapping an input function $u$ to an output function $v$. In scientific applications, this operator frequently corresponds to the evolution operator of a physical system, such as the mapping from initial conditions to future states governed by partial differential equations. One of the most influential developments in this area is the Fourier Neural Operator (FNO), which introduced spectral learning mechanisms for approximating nonlinear operators. By performing transformations in Fourier space, FNOs efficiently capture global spatial correlations and have demonstrated strong performance on a variety of partial differential equation benchmarks, including fluid dynamics, turbulence modeling, and multiphysics simulations. The success of FNOs originates from their ability to learn resolution-independent representations of physical fields. Unlike traditional convolutional architectures, which are typically tied to a fixed discretization, neural operators can generalize across different spatial resolutions and computational grids. This property makes them particularly attractive for surrogate modeling of computationally expensive numerical simulations. Another important direction is Deep Operator Network (DeepONet), which introduced a branch-trunk architecture designed to approximate nonlinear operators using separate representations of input functions and evaluation coordinates. DeepONet demonstrated the ability to learn complex mappings between functional spaces and has been applied to differential equations, inverse problems, and scientific forecasting tasks. Subsequent developments have expanded the neural operator paradigm through graph-based operators, attention mechanisms, multi-scale representations, and hybrid physics-informed formulations. These approaches have improved adaptability to irregular domains and complex geometries, extending operator learning beyond traditional grid-based formulations. Despite these advances, several limitations remain when applying neural operators to long-horizon dynamical prediction. First, many neural operator architectures primarily focus on learning global input-output mappings rather than explicitly modeling the internal mechanisms of temporal evolution. Although such models can approximate solution operators accurately within the training distribution, maintaining stability during repeated autoregressive forecasting remains challenging. Second, neural operators often rely on large collections of high-quality simulation data. In many scientific domains, generating such datasets requires expensive numerical simulations, limiting scalability and reducing applicability in situations where observations or simulations are sparse. Third, while some physics-informed neural operators incorporate differential equation constraints or conservation principles, physical knowledge is commonly introduced as an additional training objective rather than as a mechanism controlling the internal evolution of learned representations. Consequently, the learned operator may reproduce observed solution manifolds while lacking explicit guarantees of physically consistent long-term behavior. A further challenge concerns complex spatial structures. Many physical systems involve irregular meshes, evolving topologies, localized interactions, or heterogeneous components. Although recent graph neural operators address some of these limitations, most operator-learning approaches still rely on learning transformations between function spaces without explicitly representing physical interaction networks. From this perspective, graph-based physics-guided architectures provide a complementary approach. Rather than learning only the global mapping between functional spaces, PHYSGEN models physical evolution as a sequence of locally interacting graph processes guided by domain knowledge. The proposed framework combines the relational advantages of graph representations with explicit physical guidance mechanisms, enabling adaptive modeling of complex dynamical systems while maintaining long-horizon stability. Therefore, PHYSGEN should be viewed not as a replacement for neural operators, but as an alternative scientific learning paradigm that emphasizes physically guided temporal evolution, graph-based interaction modeling, and stability-aware prediction. \] 

